\section{进动、TiP与声子语言的边界}\label{sec:3-4}

在第\ref{sec:3-2}节与第\ref{sec:3-3}节中，我们建立了纯三轴转子及单准粒子耦合转子绕稳定主轴摆动运动的局域展开模型，推导了摆动特征频率在横向与纵向几何构型下的演化机制。在非轴对称原子核的转动谱学中，围绕低激发伴带的动力学本质长期存在着术语内涵与模型机制的多重竞争：一方面，唯象声子图景常将激发带直接赋予“单声子摆动（$n_w=1$）”或“双声子摆动（$n_w=2$）”标签，但在有限自旋的介观核体系中，简谐声子近似的有效性边界、波函数的空间弥散度以及内禀自由度的重取向效应需要经受严格的理论审视；另一方面，在正常形变核区（如 $A \approx 130$ 核区），伴带与基带之间的能级间距随自旋演化的特征究竟应当被理解为几何各向异性主导的横向摆动，还是归因于三维空间转轴倾斜进动（Tilted Precession, TiP）、多准粒子组态交叉、普通旋称伙伴带分裂或 $\gamma$ 软核的大振幅形状涨落，成为了近二十年低能核物理领域的重大前沿争议。

本节旨在系统阐明轴对称高 $K$ 能带进动的微观机制，严格解耦固定自旋局域旋转系振动频率与固定内禀投影跨自旋实验室能量间距；深入阐述倾斜进动（TiP）模型的粒子—转子取向几何，将有限自旋量子波函数的节点与概率分布同经典相空间轨道清晰区分；进而划定简谐声子语言在有限核中的适用界限与多重替代机制的条件竞争；最后确立结合多维电磁跃迁网络进行实验鉴别与联合判据的三轴谱学分析准则。

\subsection{轴对称高 $K$ 能带进动与转动态能量间距}\label{sec:3-4-1}

在探讨非轴对称核取向动力学之前，必须首先建立轴对称变形核中高 $K$ 态的严格物理基准。在微观多体物理中，连续均匀的轴对称量子介质（如偶偶核基态 $K=0$）在量子力学上不存在绕对称轴的独立集体转动自由度——绕空间对称轴的有限旋转对应于系统的平庸规范变换，并不产生可分辨的物理激发态（即集体转动惯量 $\mathcal{J}_\parallel = 0$）。因此，量子轴对称核不能等同于具有任意有限 $\mathcal{J}_\parallel$ 的经典刚体。

当核内未配对价准粒子产生强烈的内禀空间极化时，设未配对核子的总内禀角动量算符为 $\hat{\mathbf{J}}_{\text{intr}} = \sum_i (\hat{\mathbf{l}}_i + \hat{\mathbf{s}}_i)$。定义无量纲角动量算符 $\hat{\boldsymbol{\mathcal{I}}} \equiv \hat{\mathbf{I}}/\hbar$ 与 $\hat{\mathbf{j}}_{\text{intr}} \equiv \hat{\mathbf{J}}_{\text{intr}}/\hbar$，二者具有一致的物理量纲规范。当其在内禀对称轴（设为3轴）上的投影达到较大数值 $K \gg 1$ 时，形成高 $K$ 态。此时非零的轴向角动量分量并非源于质量分布绕对称轴的集体旋转，而是由微观费米子波包锁定所赋予的守恒内禀角动量偏置（$\hat{j}_{\text{intr}, 3} = K$）。

在转子—粒子强耦合绝热极限下，体系在垂直于对称轴方向具有集体主转动惯量 $\mathcal{J}_\perp > 0$。集体转动动能仅由垂直于对称轴的集体角动量 $\hat{\mathbf{R}}_\perp = \hbar (\hat{\boldsymbol{\mathcal{I}}}_\perp - \hat{\mathbf{j}}_{\text{intr},\perp})$ 支撑：
\begin{equation}\label{eq:3-4-high-K-H}
\hat{H}_{\text{coll}} = \frac{\hbar^2}{2\mathcal{J}_\perp} \hat{\mathbf{R}}_\perp^2 = \frac{\hbar^2}{2\mathcal{J}_\perp} \left( \hat{\boldsymbol{\mathcal{I}}}_\perp - \hat{\mathbf{j}}_{\text{intr},\perp} \right)^2.
\end{equation}
在绝热强耦合近似下，内禀态处于确定的高 $K$ 本征态，横向科里奥利退耦项 $\langle K | \hat{j}_{\text{intr},\pm} | K \rangle$ 在 $K > 1/2$ 时严格为零，横向单粒子涨落项并入内禀能 $E_{\text{intr}}$。于是，该强耦合模型下实验室系高 $K$ 态的转动能谱给出：
\begin{equation}\label{eq:3-4-high-K-spectrum}
E_K(I) - E_{\text{intr}} = \frac{\hbar^2}{2\mathcal{J}_\perp} \left[ I(I+1) - K^2 \right] \quad (I \ge K).
\end{equation}
在上述理想强耦合模型与允许的旋称/宇称选择定则下，转动带的理论带头自旋满足 $I_{\text{bandhead}} = K$。需要明确指出，这一带头恒等关系是该特定模型的代数结论，在实际在束 $\gamma$ 谱学观测中，由于向低能邻近能级的电磁退激分支、实验探测阈值或科里奥利混合对低自旋能级的微观下移，实验观测到的带头自旋不必然严格等于内禀投影量子数 $K$。

考察该转动带内相邻自旋能级在实验室系中的跃迁能量间距：
\begin{equation}\label{eq:3-4-lab-spacing}
\Delta E_{\text{lab}}(I \to I-1) \equiv E_K(I) - E_K(I-1) = \frac{\hbar^2 I}{\mathcal{J}_\perp}.
\end{equation}
在模型带头紧邻的首跨跃迁（$I = K+1 \to K$）处，式~\eqref{eq:3-4-lab-spacing} 给出首跃迁能量间距为：
\begin{equation}\label{eq:3-4-first-lab-spacing}
\Delta E_{\text{lab}}(K+1 \to K) = \frac{\hbar^2 (K+1)}{\mathcal{J}_\perp}.
\end{equation}
在大投影极限下（$K \gg 1$），该首跃迁能量间距渐近趋近于：
\begin{equation}\label{eq:3-4-asymptotic-lab-spacing}
\Delta E_{\text{lab}} \approx \frac{\hbar^2 K}{\mathcal{J}_\perp}.
\end{equation}

在此必须严格厘清局域共转参考系特征频率与实验室跨自旋能级差分的本质物理区别：
\begin{enumerate}
\item \textbf{旋转系局域频率与实验室能级差分的物理范畴解耦}：在固定自旋 $I$ 的局域随体旋转参考系中，非轴对称核围绕主轴展开的摆动特征频率（如第\ref{sec:3-2}节式~\eqref{eq:3-2-limit2-nonzero} 给出的 $\hbar \omega_{\text{wob}} = \sqrt{(\hbar I/\mathcal{J}_y - \hbar\omega_{\text{rot}})(\hbar I/\mathcal{J}_z - \hbar\omega_{\text{rot}})}$），描述的是同一自旋下转动矢量相对于内禀形状取向的小振幅振动态；而式~\eqref{eq:3-4-lab-spacing} 给出的 $\Delta E_{\text{lab}}$ 描述的是固定内禀构型下沿转动带跨越不同总角动量量子数（$I \to I-1$）的实验室能级差。二者在参考系选择、守恒量取向及实验跃迁拓扑上属于完全不同的范畴。
\item \textbf{轴对称极限下的代数联系与物理界限}：Shimizu、Matsuzaki 与 Matsuyanagi（2005）~\cite{shimizu-2005-high-k-precession} 系统探讨了微观准粒子随机相位近似（RPA）摆动形式理论的轴对称极限。当三轴核的形变在转轴（设为 $x$ 轴）周围恢复轴对称性（$\mathcal{J}_y = \mathcal{J}_z \equiv \mathcal{J}_\perp$）时，若将内禀态取为沿对称轴具有大投影 $K = \langle \Phi | \hat{J}_x | \Phi \rangle$ 的准粒子态，系统在实验室系中的进动特征频率由 $\omega_{\text{prec}} = \hbar K/\mathcal{J}_\perp^{(\text{eff})}$ 给出。在绝热 RPA 极限下，随转系微小摆动频率与实验室进动频率严格满足：
\begin{equation}\label{eq:3-4-shimizu-connection}
\omega_{\text{wob}} = \omega_{\text{prec}} - \omega_{\text{rot}},
\end{equation}
其中 $\omega_{\text{rot}} = \hbar I/\mathcal{J}_x$ 为推转参考系的旋转角速度。式~\eqref{eq:3-4-shimizu-connection} 深刻表明：在特定的轴对称绝热极限下，共转系摆动模式退化为实验室系高 $K$ 进动模减去参考系转动角频率。这一结论既揭示了三轴取向摆动与轴对称进动在特定极限下的内在代数联系，也严格划定了二者的适用界限——通用横向摆动建立在完全非对称三轴惯量张量与非平行准粒子取向之上，绝不能不加检验地将所有横向摆动与轴对称高 $K$ 进动视作平庸同构，亦不能割裂它们在特定连续极限下的动力学映射。
\end{enumerate}

\subsection{倾斜进动（TiP）模型的取向动力学描述}\label{sec:3-4-2}

针对正常形变核区（如 $A \approx 130$ 区）中伴带能量间隔随自旋反常演化的物理机制，Lawrie、Shirinda 与 Petrache（2020）基于粒子—转子模型（PRM）建立了倾斜进动（Tilted Precession, TiP）唯象取向动力学图景~\cite{lawrie-2020-tilted-precession-wobbling}。

在粒子—转子耦合体系中，系统的总哈密顿量由非轴对称核芯转动动能与价准粒子在三轴形变场中的单粒子哈密顿量共同构成。体系的三维转动态可通过总角动量球面 $|\mathbf{I}|^2 = \text{const}$ 与能量椭球面 $E = \text{const}$ 的相空间相交几何进行半经典表象：
\begin{itemize}
\item \textbf{准粒子取向与核芯转动的动力学竞争}：奇 $A$ 核中侵入的高 $j$ 准粒子（如缺中子稀土与钡镧铈区中的 $\pi h_{11/2}$ 轨道）受到强各向异性四极驱动势的作用，倾向于顺排在能量最有利的几何主轴（粒子型倾向于沿短轴 $s$，空穴型倾向于沿长轴 $l$）；而三轴形变核芯由于集体流体力学效应，倾向于绕其有效转动惯量最大的主轴（如无旋流模型下的中间轴 $m$）转动。价核子优势顺排取向与核芯惯量优势转轴的夹角竞争，导致系统的经典稳定转动取向轴脱离单一几何主轴，向主轴所张成的平面（如 $s$-$m$ 平面）发生空间倾斜。
\item \textbf{TiP 能带的物理定义与进动取向}：总角动量矢量围绕这一倾斜的稳定方向做三维空间进动。在完整粒子—转子模型（QTR）求解中，激发态对应于具有更大倾斜角或进动锥角的转动带，Lawrie 等人将这类包含完整三维自由度的转动态统称为倾斜进动（TiP）能带。
\item \textbf{摆动近似（Wobbling Approximation）的受限性质}：简谐摆动运动在数学上对应于 TiP 动力学的一个极端小振幅近似分支。只有当总角动量在另外两个垂直方向的分量平方和远小于总自旋平方（即无量纲小量参量 $f(n, I) \ll 1$），且激发能和电磁跃迁强度对激发阶数 $n$ 呈现严格的线性简谐量子化标度时，系统才严格退化为一维转动与谐振声子直接解耦的摆动图景。
\item \textbf{模型机制与唯象判据的客观界限}：准粒子取向与核芯转动的张角竞争是粒子—转子哈密顿量本征求解给出的内部模型机制，并非可以在实验室系中利用探测器沿时间流逝直接追踪的经典质点轨迹。此外，理论计算表明，倾斜进动带与基带之间的能级间距 $\Delta E(I)$ 随自旋演化的趋势（增加、下降或非单调改变）敏感地依赖于具体的单粒子轨道、核芯惯量排序及有效形变参数，能级差下降并非区分 TiP 与其他模式的排他性唯一判据；同理，亦不能脱离微观自洽输入，为主轴三轴度或自旋窗口强加跨核区通用的绝对经验界限。
\end{itemize}

\subsection{有限自旋波函数分布与经典轨迹的几何区分}\label{sec:3-4-3}

在探讨三轴核的取向激发时，极易将经典力学相空间中绘制的连续轨道切线，直观等同于微观核态的实际量子概率分布。在有限自旋的介观原子核中，必须确立严格的量子力学几何区分：
\begin{enumerate}
\item \textbf{经典相轨道的局限性}：在由能量曲面与角动量守恒面约束的经典相空间中，角动量矢量的端点沿严格确定的闭合轨线周期演化，进动锥角与倾斜取向具有明确的瞬时几何坐标。
\item \textbf{有限自旋量子波函数的离散基矢与概率分布}：真实原子核的量子本征态由包含有限维微观基底的投影波函数构成（如三维转动 $D$ 函数与准粒子态的相干叠加矢 $|\Psi_{IM}\rangle = \sum_{K,\alpha} C_{K\alpha}^I |I M K \alpha\rangle$）。该波函数在内禀空间表现为具有空间弥散度的量子概率分布，而非理想的一维闭合曲线：
\begin{itemize}
\item 在有限自旋（如中等自旋区）下，基态（$n=0$ 模式）在有效势能极小邻域呈现出局域化的量子波包。其横向角向展宽受控于局部势能曲率比与总自旋：在大自旋局域简谐近似展开下，横向角动量量子涨落满足 $\Delta \mathcal{I}_\perp \sim \sqrt{I}$，从而局部角向弥散度呈现特征性标度 $\Delta \theta \sim \Delta \mathcal{I}_\perp / I \sim 1/\sqrt{I}$。需要强调，这一标度关系严格依赖于势能面在极小点具有稳定的二次型恢复力曲率且处于远离软模的低阶激发（$n \ll I$），不能视作在所有自旋和所有核态下均普适适用的半高全宽（FWHM）；
\item 对于第一激发态（$n=1$ 模式），在理想的一维对称简谐势阱模型 $\psi_1(x) \propto x e^{-x^2/2}$ 中，其波函数在平衡极小点严格具有一个中心对称节线（节点处概率密度为零），并在两侧形成对称双峰。然而在真实的核多体哈密顿量中，由于两横向主轴惯量不对称、科里奥利相互作用的非对角耦合以及内禀宇称/旋称成分的混合，一维简谐反射对称性被普遍破坏，真实的 $n=1$ 量子态并不保证在经典转轴中心处严格具有零概率点，双峰结构亦会发生不对称扭曲，其弥散范围取决于实际势阱的约束深度，不能随意夸大为必然横跨整个半球；
\item 当体系激发能量接近相空间双曲分离线（separatrix）或局部鞍点时，非简谐展开高阶项占优，量子隧穿效应与波函数混合急剧增强，局域简谐展开失效，必须依赖完整三维对角化或投影计算。
\end{itemize}
\item \textbf{实验可观测量的客观映射}：在核物理实验中，探测器记录的是光子能量、时间戳与角度符合计数，并不存在直接对空间量子概率云密度或瞬时经典轨道进行直接成像的测量手段。理论计算求得的量子态密度仅能作为模型解释的微观媒介，对核结构取向动力学的科学判定必须完全建立在可实验验证的跃迁能谱、分支比、多极混合比及约化电磁跃迁几率 $B(E2)$ 和 $B(M1)$ 之上。
\end{enumerate}

\subsection{谐振声子语言的适用界限与替代机制竞争}\label{sec:3-4-4}

基于前述取向动力学与波函数特性的分析，我们可以严格划定简谐声子语言在核谱学中的适用边界与多重竞争机制。

\subsubsection{声子标记的局域有效性与破坏条件}

在微观粒子—转子模型或宏观转子展开中，将激发带标记为 $n_w = 1$ 或 $n_w = 2$ 摆动声子带，在动力学上仅在满足以下局域有效条件时才构成合理的物理近似：
\begin{itemize}
\item \textbf{有效谐振势阱展开条件}：系统在选定的转动取向参考轴邻域内，总有效势能面具有清晰的局域极小，且量子波包的主体被紧密约束在二次型势阱展开的有效作用半径之内，使得高阶非简谐微扰项（如四次方拉伸项与非对角反常耦合）相对于谐振能标 $\hbar \omega_{\text{wob}}$ 属于小量；
\item \textbf{有效惯量与顺排稳定性条件}：局域有效恢复力由惯量张量与内禀取向共同支撑。需要澄清的是，小振幅声子展开并不要求三个主转动惯量必须显著互异（如前文所论证的轴对称极限 $\mathcal{J}_y = \mathcal{J}_z$ 同样存在严格自洽的局域 RPA 声子展开），亦不要求价准粒子角动量 $\mathbf{j}$ 在旋转全过程中绝对刚性冻结而不发生任何微观取向调整；其核心物理要求在于，随转速演化引起的准粒子自洽重排速率远慢于横向摆动的集体振荡周期。
\end{itemize}

当体系处于低自旋区（$I \sim j$）或遭遇多准粒子带交叉区时，强烈的科里奥利力与准粒子退耦使得势能曲面变得极其平坦甚至发生极值点漂移，非简谐非对角相互作用急剧增强。然而，带交叉或低自旋扰动的出现并不意味着在所有物理情景下都会导致电磁跃迁强度的彻底淬灭或声子物理图像的瞬间瓦解；实验与微观计算应当针对具体的能级拓扑与波函数成分，实事求是地评估声子概念的适用程度及其偏差。

\subsubsection{谱学中的多重竞争机制}

在在束 $\gamma$ 谱学中，当实验测得一条伴带具有密集的 $\Delta I = 1$ 退激并与基带保持较小能量差时，除横向摆动之外，以下替代物理机制构成了强有力的条件竞争解释：
\begin{enumerate}
\item \textbf{普通旋称伙伴带（Signature Partner Bands）}：在单奇 $A$ 形变核中，同一高 $j$ 准粒子轨道（如 $\pi h_{11/2}$）分裂出的两个旋称分支（$\alpha = \pm 1/2$）在适度形变与费米面位置下，其旋称分裂能量差随自旋升高可以表现出平缓缩小或交替振荡的行为，在特定自旋区间内能够模仿出伴带能级差随自旋下降的表象；
\item \textbf{相邻准粒子组态交叉（Band Crossing）}：奇 $A$ 核在高自旋转动过程中，极易经历破对（例如一对高 $j$ 中子对或质子对的顺排）引发的多准粒子组态交叉。在能级交叉区域，由于两组态之间的残余相互作用发生能级排斥与波函数混合，往往在有限自旋跨度内诱发伴带能谱的局部压缩甚至近简并；
\item \textbf{动态软核核芯耦合（Soft-core Vibrational Coupling）}：对于广泛分布于 $A \approx 130$ 核区的弱变形核，其势能面对 $\gamma$ 自由度展现出显著的柔软性（$\gamma$-softness）。在此类体系中，核芯低激发谱包含显著的 $\gamma$ 振动与大振幅四极形状涨落。未配对准粒子与软核芯表面振动的动力学耦合，在不需要假定深陷刚性三轴极小的前提下，同样能够重现实验观测到的低激发带结构及其电磁退激特征（Petrache 等人 2006，Lv 等人 2022，Nomura 2022）~\cite{petrache-2006-near-degenerate-chiral-misinterpretation,lv-2022-evidence-against-wobbling-135pr,nomura-2022-questioning-wobbling-ibfm}。
\end{enumerate}

在实际核谱学研究中，必须严格恪守实验数据的测量边界与不确定度传播规律。对于处于前沿争议中的候选核素，若相关实验测量仅确立了激发能级纲图而缺乏精确的能级平均寿命（如皮秒量级的反冲距离法或多普勒位移衰减法测量），则该项实验数据无法独立提供带间跃迁的绝对电磁几率 $B(E2)$ 与 $B(M1)$。在数据不完备的情形下，研究必须明确标注测量缺口与物理边界，坚决避免在缺乏直接寿命测量的支撑下虚构或武断推断绝对跃迁强度。

\subsection{实验观测量裁决能力与三轴谱学分析准则}\label{sec:3-4-5}

面对上述多重竞争模型的挑战，单一观测量（如孤立的能谱能级差或定性的跃迁存在性）无法对三轴取向动力学做出排他性裁决。现代核结构谱学确立了联合多维实验可观测量进行综合判决的分析准则：
\begin{enumerate}
\item \textbf{能谱间距与激发趋势（$\Delta E(I)$）}：提供动力学模式的初步判据，建立伴带相对于基带在整个可测自旋区间的演化轮廓，用以初筛是否存在旋称反转、带交叉或特征频率分叉；
\item \textbf{方向角关联（DCO 比值）与线偏振度（Linear Polarization）}：用于测定级联退激的跃迁多极度（区分电四极 $E2$ 与磁偶极 $M1$）及多极混合比 $\delta(E2/M1)$。定向核角分布与角关联理论表明（Lange、Kumar 与 Hamilton 1982）~\cite{lange-kumar-hamilton-1982-multipole-admixtures}，方向角关联比值 $R_{\text{DCO}}$ 关于多极混合比 $\delta$ 表现为二次有理分式方程。求解实验方程 $R_{\text{DCO}}(\delta) = R_{\text{exp}}$ 时：
\begin{itemize}
\item 在代数上一般对应于两个实数根、重根或无实解（当实验实测值因统计涨落或未计入的退取向效应超出几何极限极大/极小值时），并不保证恒有“一显著大根与一显著小根”；
\item 测量 Compton 散射线偏振度 $P$ 能够提供独立的电多极与磁多极宇称鉴别，但在偏振探测灵敏度有限或两根所对应的理论偏振值处于实验误差包络内时，测量 $P$ 并不保证能彻底消除多解歧义性；
\item 多极混合比 $\delta(E2/M1)$ 在物理上严格定义为具有完整电磁算符量纲因子（$M1$ 采用核磁子 $\mu_N$，$E2$ 采用 $e\cdot\mathrm{fm}^2$，并包含跃迁光子能量因子 $E_\gamma/\hbar c$）的无量纲辐射发射几率振幅之比（与第\ref{sec:5-9}节式~\eqref{eq:5-9-delta-def} 保持一致规范），不能简单混同于无量纲裸约化矩阵元之比；
\end{itemize}
\item \textbf{能级寿命测量与绝对跃迁几率（$B(E2)$ 和 $B(M1)$）}：利用多普勒位移衰减法（DSAM）或反冲距离法（RDM）测出激发态平均寿命 $\tau$，结合光子分枝比、内转换电子系数及严格反演的多极混合比 $\delta$，提取绝对跃迁几率 $B(E2)_{\text{inter}}$ 与 $B(M1)_{\text{inter}}$。在理想的纯转子与冻结顺排小振幅摆动模型极限下（Bohr 与 Mottelson 1975，Frauendorf 与 D{\"o}nau 2014）~\cite{book-bohr-mottelson-nuclear-structure-vol2,frauendorf-2014-transverse-wobbling}，带间跃迁表现出增强的集体 $E2$ 强度，且在领先阶呈现 $B(E2)_{\text{inter}} \sim 1/I$ 的渐近下降趋势，同时伴随较小的 $M1$ 跃迁强度；然而必须清醒认识到，$1/I$ 标度与强 $M1$ 抑制是特定理想局部模型在大自旋展开下的极限结论，在真实复杂核体系中，组态重构与微观极化会导致其实验值发生系统性偏离，不能当作普适绝对定律；
\item \textbf{微观与唯象理论模型的前向联合预言检验}：最终的物理裁决必须依靠自洽微观密度泛函、全三维准粒子转子模型（QTR）或多轴推转微观投影框架，在统一且自洽的参数输入下，对能级能谱、寿命衰变链、混合比以及电磁跃迁强度进行前向联合预言，并与竞争性物理机制（旋称分裂、组态交叉、软核涨落）进行定量对账。
\end{enumerate}

对非轴对称核转动态的物理确认，本质上是在不断提升的多维实验精密度与自洽理论约束下，对竞争性物理图像进行的严密证伪与综合鉴别过程。
